\newpage
\section{Oszilloskopbilder}
Um die Auswertung nicht noch mehr mit Bilder voll zu spammen, werden manche erst hier eingefügt.
\subsection*{\RNum{1} Halbwertszeit für Konfiguration 1}
\xfigure{htbp}{width=0.8\textwidth}{A1_Konf1.png}{Halbwertszeit mit Cursorn bestimmen}{Halbwertszeit mit Cursorn bestimmen}

\subsection*{\RNum{3} Bestimmen der Grenzfrequenzen mit der 3dB Methode}
\xfigure{htbp}{width=0.8\textwidth}{A3_Tiefpass_Frqugang.png}{Grenzfrequenz mit 3dB Methode}{Grenzfrequenz des Tiefpasses wird durch einen Amplitudenabfall von 3dB bestimmt}
Diese Methode wurde analog auf den Hochpass angewandt. Die untersuchten Bilder, jedoch ohne die ikonische Cursorposition sind in den folgenden Abbildungen zu finden. 

\begin{figure}[htbp]
    {\captionsetup{list=false}
    \begin{subfigure}{0.48\textwidth}
        \centering
        \includegraphics[width=\linewidth]{A3_hochfpass_Phaseplot.png}
        \caption{Hochpasses (CR-Glied)}
        % \label{fig:3.1}
    \end{subfigure}
    \hfill
    \begin{subfigure}{0.48\textwidth}
        \centering
        \includegraphics[width=\linewidth]{A3_Tiefpass_Phaseplot.png}
        \caption{Tiefpass (RC-Glied)}
        % \label{fig:3.2}
    \end{subfigure}}
    \caption[Hoch- und Tiefpass Frequenzgang am Oszilloskop]{Hoch- und Tiefpass Frequenzgang, bei \(C=\qty{47}{\nano\farad}\) und \(R=\qty{1}{\kilo\ohm}\)}
      \label{fig:freqgängenoszi}
\end{figure}
\FloatBarrier

\subsection*{\RNum{9} ungefiltertes Eingangssignal}

\xfigure{htbp}{width=0.9\textwidth}{A8_signalforms_grundsignal.png}{Eingangssignal ungefiltert}{Eingangssignal, ungefiltert}
\clearpage

\section{Memes}
\xfigure{htbp}{width=0.9\textwidth}{Sunsetboot.png}{\emph{Meme:} Sonnenuntergang}{Sonnenuntergang von der Ernst-Walz-Brücke aus}


\section{Code-Anhang}
Bei den Funktionen habe ich teilweise Dinge von verschiedenen Altprotokollen zusammengesucht, und auf meine Bedürfnisse umgeschrieben. Die konkreten Berechnung sind dann eigenständig.
\includepdf[pages=-]{../Code/241/241.pdf}
